% page 356 of the facsimile (image p-355 of the scan)
% block 154.9 x 246.3 mm ; left margin 8.0 mm
\pgc{-12.700mm}{0.000mm}{
\l{\marge[79.046mm]{26.280mm}{\ornamento{14.228mm}{ornaments/letters/letter-M.png}}\cel{3}348}
\l{}
\l{}
\l{}
\l{}
\l{}
\l{}
\l{}
\l{}
\l{}
\l{ma. Konjunciono qua indikas la interopozeso di la du pro-}
\l{\cel{3}pozicioni quin lu ligas, unionas : (1) per restriktar}
\l{\cel{3}to quo jus dicesis ; (2) per expresar ideo qua kontras-}
\l{\cel{3}tas kun olta qua jus expresesis ; (3) per objecionar a}
\l{\cel{3}to quo jus dicesis. - FIS.}
\l{}
\l{\sou{macedonio.}\cel{1}(\sou{koquarto}). Disho en qua legumi diversa esas pre-}
\l{\cel{3}parita kune. Entremeso sukroza en qua frukti diversa esas}
\l{\cel{3}asemblita aden quaza jeleo. - EFIS.}
\l{}
\l{\sou{macerar}. (\sou{trans. e netrans.})\cel{2}Trempar (substanco) aden liqui-}
\l{\cel{3}do kolda, e lasar (ta substanco) en ta liquido, por depre-}
\l{\cel{3}nar la solveblaji. - DEFIS.}
\l{}
\l{\sou{\textquotedbl{}macfarlane\textquotedbl{}}. Mantelo sen maniki, kun aperturi tra qui pa-}
\l{\cel{3}sas la brakii, sur qui kovre-pendas pelerino qua finas ye}
\l{\cel{3}la tayo di la persono. -}
\l{}
\l{\sou{\textquotedbl{}mackintosh}\textquotedbl{}. Sorto di muliero-mantelo di qua la stofo esas}
\l{\cel{3}nepermeebla.}
\l{}
\l{\sou{\textquotedbl{}madapolam\textquotedbl{}.}\cel{2}Texuro de kotono blanka, di qua la grano esas}
\l{\cel{3}tre salianta e tre apretita - meze inter kaliko e perkalo,}
\l{\cel{3}tre uzata por facar la avantali koqueyala.}
\l{}
\l{\sou{madeleno}. Mikra kuko di qua la pasto esas kompakta. DEFIS.}
\l{}
\l{\sou{madono}.\cel{2}Statuo, pikturo, qua reprezentas (che la kristani)}
\l{\cel{3}la virgino Maria. - DEFIRS.}
\l{}
\l{\sou{\textquotedbl{}madra}s\textquotedbl{}Kapovesto ek fulardo de stofo di qua la fili warpa}
\l{\cel{3}esas ek silko, e la wefto ek kotono - e qua, origine,}
\l{\cel{3}fabrikesis en la urbo Madras (India).}
\l{}
\l{\sou{madreporo}. Genero de polipi agregita, asemblajo de celuli}
\l{\cel{3}kalka qui interkomunikas, e kreskas til formacar strati,}
\l{\cel{3}rifi, insuleti, e c. - L.\cel{1}\sou{madrepora}. DEFIS.}
\l{}
\l{\sou{madrigalo}. Kurta versajo, maxim-multa-kaze turnita a damo,}
\l{\cel{3}e qua expresas penso tenera o galanta. - DEFIRS.}
\l{}
\l{\sou{maestro}. Eminento pri arto o cienco. - DEFIRS.}
\l{}
\l{\sou{mago.}I.Sacerdoto di la religio di la Persi antiqua, adoranti}
\l{\cel{3}di fairo. - II. (\sou{biblo)}. Un ek le tri personegi - quin la}
\l{\cel{3}tradiciono transformis a reji - qui venis adorar Iesu en}
\l{\cel{3}lua kripo, en urbo Betleem. - DEFIS.}
}
